% ===============================================================
% Document Class ================================================
\documentclass{article}

% ===============================================================
% Graphic Packages ==============================================
\usepackage{graphicx}       % Permite a inclusão de imagens no documento.
\usepackage{tocloft} % Pacote para customização de listas de conteúdos, figuras e tabelas
\usepackage{tikz}           % Ferramenta poderosa para criar gráficos programaticamente dentro do LaTeX.
\usetikzlibrary{calc}       % Extensão da biblioteca TikZ que permite cálculos mais complexos de coordenadas.
\usepackage[portuguese]{babel}
\usepackage{tocloft} % Pacote para customização de listas
\renewcommand{\listfigurename}{Lista de Figuras} % Altera o título para português
\AtBeginDocument{\renewcommand{\contentsname}{Sumário}} % Altera nome de conteudo para sumário
% ===============================================================
% Mathematical Tools ============================================
\usepackage{amsmath}        % Melhora a aparência e a flexibilidade de comandos matemáticos.
\usepackage{siunitx}        % Facilita o uso de unidades do Sistema Internacional e ajuda a formatar números complexos.

% ===============================================================
% Font and Text Appearance ======================================
\usepackage{mathptmx}       % Altera a fonte padrão do documento para Times New Roman.

% ===============================================================
% Table of Contents Customization ===============================
\usepackage{tocloft}  % Oferece controle total sobre a aparência das listas de conteúdos, figuras, tabelas, etc.

% ===============================================================
% Figure Positioning ============================================
\usepackage{float}     % Melhora a interface para definir o posicionamento de objetos flutuantes como figuras e tabelas.

% ===============================================================
% Paragraph Spacing and Indentation =============================
\usepackage{setspace}  % Permite o ajuste fino do espaçamento entre linhas.
\usepackage{indentfirst} % Adiciona indentação ao primeiro parágrafo de cada seção.

% ===============================================================
% Page Layout ===================================================
\usepackage[a4paper, top=3cm, bottom=2cm, left=3cm, right=2cm]{geometry}  % Define as margens de todo o documento.
\setlength{\parindent}{4em}  % Define o tamanho da indentação para todos os parágrafos.
\setlength{\emergencystretch}{3em}

% ===============================================================
% Section Heading Customization =================================
\makeatletter
\renewcommand\paragraph{\@startsection{paragraph}{4}{\z@}%
    {2ex plus 1ex minus .2ex}%
    {1em}%
    {\normalfont\normalsize\bfseries}}
\makeatother

% ===============================================================
% Urls ==========================================================
\usepackage{hyperref} % Para criar links clicáveis
\hypersetup{
    colorlinks=true,
    linkcolor=blue,
    filecolor=magenta,
    urlcolor=blue,
    citecolor=blue,
    pdfborder={0 0 0},  % Remove o quadrado ao redor dos links
    breaklinks=true  % Permite que os links sejam quebrados entre as linhas
}
\usepackage{xurl} % Melhora a quebra de links em qualquer lugar

% ===============================================================
% Bloco de código ===============================================
\usepackage{listings}
\usepackage{xcolor}
\usepackage{courier}
\lstset{
  backgroundcolor=\color[RGB]{249,246,239},
  basicstyle=\ttfamily\footnotesize,
  breaklines=true,
  frame=single,
  numbers=left,
  numberstyle=\tiny\color{gray},
  keywordstyle=\color[RGB]{40,40,255},
  commentstyle=\color[RGB]{0,125,0},
  stringstyle=\color[RGB]{255,0,0},
  showstringspaces=false,
  rulecolor=\color{black},
  captionpos=b,
  abovecaptionskip=5pt,
  belowcaptionskip=5pt,
  xleftmargin=0.15\textwidth,
  xrightmargin=0.15\textwidth,
  morecomment=[l]{//}
}
\renewcommand{\lstlistlistingname}{Lista de Códigos} % Comando para renomear o título da Lista de Códigos

% ===============================================================
% Begin Document ================================================
\begin{document}

% ===============================================================
% Capa ==========================================================
\begin{titlepage}
    \centering
    % Desenhar a margem
    \begin{tikzpicture}[remember picture, overlay]
        \draw[line width = 4pt] ($(current page.north west) + (20mm, -20mm)$) rectangle ($(current page.south east) + (-10mm,10mm)$);
        \draw[line width = 1pt] ($(current page.north west) + (21.5mm,-21.5mm)$) rectangle ($(current page.south east) + (-11.5mm,11.5mm)$);
    \end{tikzpicture}

    % Linha de logos
    \includegraphics[height=0.151\textwidth]{header/contra-capa/assets/uerj.png}\hfill
    \includegraphics[height=0.15\textwidth]{header/contra-capa/assets/iprj.jpeg}\hfill

    \vspace{2cm} % Espaço vertical após os logos

    % Informações do curso
    {\Large\bfseries Instituto Politécnico do Estado do Rio de Janeiro \par}
    \vspace{0.5cm}
    {\large Curso de Engenharia da Computação \par}

    \vspace{3cm} % Espaço vertical antes do título da atividade

    % Nome do estudante
    {\large\bfseries Guilherme Cagide Fialho \par}

    \vspace{3cm}

    % Título da atividade
    {\large\bfseries Trabalho 2 - Microsoft\par}

    \vfill % Empurra o restante para o fundo da página

    % Local e ano
    {\large\bfseries Nova Friburgo \par}
    \vspace{0.3cm}
    {\large\bfseries 2024 \par}
\end{titlepage}
\newpage % Começa uma nova página após a capa


% ===============================================================
% Contra Capa ===================================================
\begin{titlepage}
    \centering
    % Desenhar a margem
    \begin{tikzpicture}[remember picture, overlay]
        \draw[line width = 4pt] ($(current page.north west) + (20mm, -20mm)$) rectangle ($(current page.south east) + (-10mm,10mm)$);
        \draw[line width = 1pt] ($(current page.north west) + (21.5mm,-21.5mm)$) rectangle ($(current page.south east) + (-11.5mm,11.5mm)$);
    \end{tikzpicture}

    % Linha de logos
    \includegraphics[height=0.151\textwidth]{header/contra-capa/assets/uerj.png}\hfill
    \includegraphics[height=0.15\textwidth]{header/contra-capa/assets/iprj.jpeg}\hfill

    \vspace{2cm} % Espaço vertical após os logos

    % Informações do curso
    {\Large\bfseries Instituto Politécnico do Estado do Rio de Janeiro \par}
    \vspace{0.5cm}
    {\large Graduação em Engenharia da Computação \par}

    \vspace{3cm} % Espaço vertical antes do título da atividade

    % Nome do estudante
    {\large Guilherme Cagide Fialho \par}

    \vspace{1.5cm}

    % Título da atividade
    {\large\bfseries Trabalho 5 - Amazon \par}

    \vspace{1cm} % Espaço vertical antes do título da atividade

    % Informações do projeto
    \begin{flushright}
        \begin{minipage}{0.5\textwidth}
            \large
            \raggedleft % Garante que o texto dentro da minipage seja alinhado à direita
            Trabalho da Disciplina Estudo de Casos Empresariais
        \end{minipage}
    \end{flushright}

    \vspace{1.5cm}

    % Informações do orientador
    \begin{flushleft}
        \begin{minipage}{0.5\textwidth}
            \large
            \raggedright
            Orientador: Prof. Leandro de Amorim Ratamero
        \end{minipage}
    \end{flushleft}


    \vfill % Empurra o restante para o fundo da página

    % Local e ano
    {\large Nova Friburgo \par}
    \vspace{0.3cm}
    {\large 2024 \par}
\end{titlepage}
\newpage % Começa uma nova página após a capa


% ===============================================================
% Resumo ========================================================
% \begin{titlepage}
    \thispagestyle{empty} % Remove números de página
    \setstretch{1.5} % Espaçamento entre linhas, certifique-se de que o pacote setspace está incluído em document.tex

    \begin{center}
        \textbf{\Large RESUMO}
    \end{center}


    \vspace{1cm} % Espaço vertical

    \noindent CAGIDE FIALHO, G. Relatório do projeto de Modelagem
    e controle de sistemas. 2024. 54 f. Trabalho de Conclusão de Disciplina Modelagem e Controle de Sistemas (Graduação em
    Engenharia da computação) – Graduação em Engenharia da Computação, Universidade
    do Estado do Rio de Janeiro, Nova Friburgo, 2024.

    \vspace{0.4cm} % Espaço vertical

    Este trabalho explora a modelagem e o controle de um sistema dinâmico do tipo massa-mola-amortecedor, utilizando a plataforma Scilab para desenvolvimento e simulação. O foco do estudo está na implementação de modelos matemáticos para descrever a dinâmica do sistema e na análise de sua resposta sob diversas condições iniciais, sem a aplicação de forças externas. Utilizando a ferramenta Xcos, um componente gráfico do Scilab, realizamos simulações que permitiram uma análise visual e quantitativa das respostas transientes do sistema. O estudo destaca a influência dos parâmetros físicos, como a massa, o coeficiente de amortecimento, e a constante da mola, nas características de resposta do sistema. Além disso, técnicas de controle foram empregadas para ajustar a resposta do sistema, demonstrando como o amortecimento pode contribuir para a estabilização após perturbações e enfatizando a relevância de uma parametrização cuidadosa para alcançar um comportamento eficaz do sistema. Este projeto contribui para a compreensão das teorias de controle aplicáveis em sistemas mecânicos e outros contextos de sistemas dinâmicos na engenharia.
    \vspace{0.4cm} % Espaço vertical

    \textbf{Palavras-chave}: Modelagem e Controle, Sistema Massa-Mola-Amortecedor, Simulação, Scilab, Xcos.
\end{titlepage}


% ===============================================================
% Abstract ======================================================
% \begin{titlepage}
    \thispagestyle{empty} % Remove page numbers
    \setstretch{1.5} % Line spacing, make sure the setspace package is included in document.tex

    \begin{center}
        \textbf{\Large ABSTRACT}
    \end{center}

    \vspace{1cm} % Vertical space

    \noindent CAGIDE FIALHO, G. Report on the Modeling and Control Systems Project. 2024. 54 p. Final Project for the Modeling and Control Systems Course (Bachelor’s degree in Computer Engineering) - Bachelor’s Degree in Computer Engineering, State University of Rio de Janeiro, Nova Friburgo, 2024.

    \vspace{0.4cm} % Vertical space

    This work explores the modeling and control of a dynamic system of the mass-spring-damper type, using the Scilab platform for development and simulation. The focus of the study is on implementing mathematical models to describe the system dynamics and analyzing its response under various initial conditions without the application of external forces. Using Xcos, a graphical component of Scilab, we conducted simulations that allowed for visual and quantitative analysis of the system's transient responses. The study highlights the influence of physical parameters, such as mass, damping coefficient, and spring constant, on the system's response characteristics. Furthermore, control techniques were employed to adjust the system's response, demonstrating how damping can contribute to stabilization after disturbances and emphasizing the importance of careful parameterization to achieve effective system behavior. This project contributes to the understanding of control theories applicable to mechanical systems and other dynamic system contexts in engineering.
    \vspace{0.4cm} % Vertical space

    \textbf{Keywords}: Modeling and Control, Mass-Spring-Damper System, Simulation, Scilab, Xcos.
\end{titlepage}


% ===============================================================
% Lista de figuras ==============================================
% \newpage
\setstretch{1.5} % Espaçamento entre linhas, certifique-se de que o pacote setspace está incluído em document.tex

% % Estilo de citação (plain, abbrv, ieeetr, etc.)
% \bibliographystyle{IEEEtran}  % Estilo que inclui URLs
% % \bibliographystyle{plain}

% \bibliography{
%     perguntas/1-foundation/latex/references,
%     perguntas/2-graduation/latex/references,
%     perguntas/3-main-productions/latex/references,
%     perguntas/4-location/latex/references,
%     perguntas/5-competitors/latex/references,
%     perguntas/6-protection/latex/references
% }

% Estilo de citação (plain, abbrv, ieeetr, etc.)
\bibliographystyle{IEEEtran}  % Estilo que inclui URLs
% \bibliographystyle{plain}

\bibliography{perguntas/1-foundation/latex/references,perguntas/2-graduation/latex/references,perguntas/3-main-products/latex/references.bib,perguntas/4-location/latex/references,perguntas/5-competitors/latex/references,perguntas/6-protection/latex/references, perguntas/7-profit/latex/references, perguntas/8-train/latex/references, perguntas/9-social-programs/latex/references, perguntas/10-green-purposes/latex/references}


% ===============================================================
% Lista de códigos ==============================================
% \newpage
\setstretch{1.5} % Espaçamento entre linhas, certifique-se de que o pacote setspace está incluído em document.tex

% % Estilo de citação (plain, abbrv, ieeetr, etc.)
% \bibliographystyle{IEEEtran}  % Estilo que inclui URLs
% % \bibliographystyle{plain}

% \bibliography{
%     perguntas/1-foundation/latex/references,
%     perguntas/2-graduation/latex/references,
%     perguntas/3-main-productions/latex/references,
%     perguntas/4-location/latex/references,
%     perguntas/5-competitors/latex/references,
%     perguntas/6-protection/latex/references
% }

% Estilo de citação (plain, abbrv, ieeetr, etc.)
\bibliographystyle{IEEEtran}  % Estilo que inclui URLs
% \bibliographystyle{plain}

\bibliography{perguntas/1-foundation/latex/references,perguntas/2-graduation/latex/references,perguntas/3-main-products/latex/references.bib,perguntas/4-location/latex/references,perguntas/5-competitors/latex/references,perguntas/6-protection/latex/references, perguntas/7-profit/latex/references, perguntas/8-train/latex/references, perguntas/9-social-programs/latex/references, perguntas/10-green-purposes/latex/references}


% ===============================================================
% Lista de siglas e abreviaturas ================================
% \newpage
\setstretch{1.5} % Espaçamento entre linhas, certifique-se de que o pacote setspace está incluído em document.tex

% % Estilo de citação (plain, abbrv, ieeetr, etc.)
% \bibliographystyle{IEEEtran}  % Estilo que inclui URLs
% % \bibliographystyle{plain}

% \bibliography{
%     perguntas/1-foundation/latex/references,
%     perguntas/2-graduation/latex/references,
%     perguntas/3-main-productions/latex/references,
%     perguntas/4-location/latex/references,
%     perguntas/5-competitors/latex/references,
%     perguntas/6-protection/latex/references
% }

% Estilo de citação (plain, abbrv, ieeetr, etc.)
\bibliographystyle{IEEEtran}  % Estilo que inclui URLs
% \bibliographystyle{plain}

\bibliography{perguntas/1-foundation/latex/references,perguntas/2-graduation/latex/references,perguntas/3-main-products/latex/references.bib,perguntas/4-location/latex/references,perguntas/5-competitors/latex/references,perguntas/6-protection/latex/references, perguntas/7-profit/latex/references, perguntas/8-train/latex/references, perguntas/9-social-programs/latex/references, perguntas/10-green-purposes/latex/references}


% ===============================================================
% Materiais utilizados ==========================================
% \newpage
\setstretch{1.5} % Espaçamento entre linhas, certifique-se de que o pacote setspace está incluído em document.tex

% % Estilo de citação (plain, abbrv, ieeetr, etc.)
% \bibliographystyle{IEEEtran}  % Estilo que inclui URLs
% % \bibliographystyle{plain}

% \bibliography{
%     perguntas/1-foundation/latex/references,
%     perguntas/2-graduation/latex/references,
%     perguntas/3-main-productions/latex/references,
%     perguntas/4-location/latex/references,
%     perguntas/5-competitors/latex/references,
%     perguntas/6-protection/latex/references
% }

% Estilo de citação (plain, abbrv, ieeetr, etc.)
\bibliographystyle{IEEEtran}  % Estilo que inclui URLs
% \bibliographystyle{plain}

\bibliography{perguntas/1-foundation/latex/references,perguntas/2-graduation/latex/references,perguntas/3-main-products/latex/references.bib,perguntas/4-location/latex/references,perguntas/5-competitors/latex/references,perguntas/6-protection/latex/references, perguntas/7-profit/latex/references, perguntas/8-train/latex/references, perguntas/9-social-programs/latex/references, perguntas/10-green-purposes/latex/references}


% ===============================================================
% Sumário =======================================================
% \newpage
\setstretch{1.5} % Espaçamento entre linhas, certifique-se de que o pacote setspace está incluído em document.tex

% % Estilo de citação (plain, abbrv, ieeetr, etc.)
% \bibliographystyle{IEEEtran}  % Estilo que inclui URLs
% % \bibliographystyle{plain}

% \bibliography{
%     perguntas/1-foundation/latex/references,
%     perguntas/2-graduation/latex/references,
%     perguntas/3-main-productions/latex/references,
%     perguntas/4-location/latex/references,
%     perguntas/5-competitors/latex/references,
%     perguntas/6-protection/latex/references
% }

% Estilo de citação (plain, abbrv, ieeetr, etc.)
\bibliographystyle{IEEEtran}  % Estilo que inclui URLs
% \bibliographystyle{plain}

\bibliography{perguntas/1-foundation/latex/references,perguntas/2-graduation/latex/references,perguntas/3-main-products/latex/references.bib,perguntas/4-location/latex/references,perguntas/5-competitors/latex/references,perguntas/6-protection/latex/references, perguntas/7-profit/latex/references, perguntas/8-train/latex/references, perguntas/9-social-programs/latex/references, perguntas/10-green-purposes/latex/references}


% ===============================================================
% Introdução ====================================================
% \newpage
\setstretch{1.5} % Espaçamento entre linhas, certifique-se de que o pacote setspace está incluído em document.tex

% % Estilo de citação (plain, abbrv, ieeetr, etc.)
% \bibliographystyle{IEEEtran}  % Estilo que inclui URLs
% % \bibliographystyle{plain}

% \bibliography{
%     perguntas/1-foundation/latex/references,
%     perguntas/2-graduation/latex/references,
%     perguntas/3-main-productions/latex/references,
%     perguntas/4-location/latex/references,
%     perguntas/5-competitors/latex/references,
%     perguntas/6-protection/latex/references
% }

% Estilo de citação (plain, abbrv, ieeetr, etc.)
\bibliographystyle{IEEEtran}  % Estilo que inclui URLs
% \bibliographystyle{plain}

\bibliography{perguntas/1-foundation/latex/references,perguntas/2-graduation/latex/references,perguntas/3-main-products/latex/references.bib,perguntas/4-location/latex/references,perguntas/5-competitors/latex/references,perguntas/6-protection/latex/references, perguntas/7-profit/latex/references, perguntas/8-train/latex/references, perguntas/9-social-programs/latex/references, perguntas/10-green-purposes/latex/references}


% ===============================================================
% Atividades ====================================================
\section{Caso: Itautec S/A}
\subsection*{Falta de Inovação e Adaptação ao Mercado}
A Itautec, conhecida por sua tecnologia e automação, falhou em inovar e adaptar-se ao mercado de TI, o que levou ao seu declínio operacional e financeiro.
\begin{itemize}
    \item \textbf{Inovação Insuficiente}: Não conseguiu acompanhar as rápidas inovações tecnológicas, perdendo espaço para concorrentes mais ágeis.
    \item \textbf{Gestão Estratégica Deficiente}: Falhas na gestão estratégica e na execução de um modelo de negócio sustentável levaram a resultados financeiros negativos.
\end{itemize}


\section{Caso: Lojas Brasileiras (Lobrás)}
\subsection*{Gestão Financeira Inadequada}
A Lobrás enfrentou falhas na gestão financeira após mudança de controle, acumulando dívidas insustentáveis.
\begin{itemize}
    \item \textbf{Endividamento Excessivo}: Acumulação de dívidas insustentáveis devido a gestão ineficiente.
    \item \textbf{Falta de Planejamento Estratégico}: Não adaptou seu modelo de negócios às mudanças do mercado varejista, resultando em perda de competitividade.
          Em 1999, as dívidas levaram a rede a fechar definitivamente as portas.
\end{itemize}



\section{Caso: Lojas Mappin}
\subsection*{Falhas na Expansão e Gestão Operacional}
O Mappin expandiu rapidamente sem uma estrutura adequada, resultando em falhas operacionais graves.
\begin{itemize}
    \item \textbf{Expansão Descontrolada}: Crescimento acelerado sem a estrutura necessária para suportar as novas operações.
    \item \textbf{Gestão Operacional Deficiente}: Falhas na administração das operações e na adaptação a novos mercados e desafios.
\end{itemize}

\section*{Conclusão}
Os casos estudados acima ilustram claramente como a falta de gestão estratégica, financeira e operacional pode levar empresas ao fracasso. Esses exemplos servem de lição para que novos empreendimentos priorizem uma inovação constante e uma gestão mais cuidadosa, fundamentais para se manterem competitivos e sustentáveis no dinâmico mercado brasileiro, garantindo a sobrevivência e o crescimento das empresas no longo prazo.

 % Completa e Validada

% ===============================================================
% Conclusão geral ===============================================
% \newpage
\setstretch{1.5} % Espaçamento entre linhas, certifique-se de que o pacote setspace está incluído em document.tex

% % Estilo de citação (plain, abbrv, ieeetr, etc.)
% \bibliographystyle{IEEEtran}  % Estilo que inclui URLs
% % \bibliographystyle{plain}

% \bibliography{
%     perguntas/1-foundation/latex/references,
%     perguntas/2-graduation/latex/references,
%     perguntas/3-main-productions/latex/references,
%     perguntas/4-location/latex/references,
%     perguntas/5-competitors/latex/references,
%     perguntas/6-protection/latex/references
% }

% Estilo de citação (plain, abbrv, ieeetr, etc.)
\bibliographystyle{IEEEtran}  % Estilo que inclui URLs
% \bibliographystyle{plain}

\bibliography{perguntas/1-foundation/latex/references,perguntas/2-graduation/latex/references,perguntas/3-main-products/latex/references.bib,perguntas/4-location/latex/references,perguntas/5-competitors/latex/references,perguntas/6-protection/latex/references, perguntas/7-profit/latex/references, perguntas/8-train/latex/references, perguntas/9-social-programs/latex/references, perguntas/10-green-purposes/latex/references}


% ===============================================================
% Referencias ===================================================
\newpage
\setstretch{1.5} % Espaçamento entre linhas, certifique-se de que o pacote setspace está incluído em document.tex

% % Estilo de citação (plain, abbrv, ieeetr, etc.)
% \bibliographystyle{IEEEtran}  % Estilo que inclui URLs
% % \bibliographystyle{plain}

% \bibliography{
%     perguntas/1-foundation/latex/references,
%     perguntas/2-graduation/latex/references,
%     perguntas/3-main-productions/latex/references,
%     perguntas/4-location/latex/references,
%     perguntas/5-competitors/latex/references,
%     perguntas/6-protection/latex/references
% }

% Estilo de citação (plain, abbrv, ieeetr, etc.)
\bibliographystyle{IEEEtran}  % Estilo que inclui URLs
% \bibliographystyle{plain}

\bibliography{perguntas/1-foundation/latex/references,perguntas/2-graduation/latex/references,perguntas/3-main-products/latex/references.bib,perguntas/4-location/latex/references,perguntas/5-competitors/latex/references,perguntas/6-protection/latex/references, perguntas/7-profit/latex/references, perguntas/8-train/latex/references, perguntas/9-social-programs/latex/references, perguntas/10-green-purposes/latex/references}


% ===============================================================
% Anexo ou Apendices ============================================
% \newpage
\setstretch{1.5} % Espaçamento entre linhas, certifique-se de que o pacote setspace está incluído em document.tex

% % Estilo de citação (plain, abbrv, ieeetr, etc.)
% \bibliographystyle{IEEEtran}  % Estilo que inclui URLs
% % \bibliographystyle{plain}

% \bibliography{
%     perguntas/1-foundation/latex/references,
%     perguntas/2-graduation/latex/references,
%     perguntas/3-main-productions/latex/references,
%     perguntas/4-location/latex/references,
%     perguntas/5-competitors/latex/references,
%     perguntas/6-protection/latex/references
% }

% Estilo de citação (plain, abbrv, ieeetr, etc.)
\bibliographystyle{IEEEtran}  % Estilo que inclui URLs
% \bibliographystyle{plain}

\bibliography{perguntas/1-foundation/latex/references,perguntas/2-graduation/latex/references,perguntas/3-main-products/latex/references.bib,perguntas/4-location/latex/references,perguntas/5-competitors/latex/references,perguntas/6-protection/latex/references, perguntas/7-profit/latex/references, perguntas/8-train/latex/references, perguntas/9-social-programs/latex/references, perguntas/10-green-purposes/latex/references}


% ===============================================================
% End Document ==================================================
\end{document}

